% =============================================================================
%  Sovereign Research Matrix: Formally Verified Computational Architecture
%  Author: Rafael D. De Paz
% =============================================================================
\documentclass[12pt, a4paper]{article}

\usepackage[utf8]{inputenc}
\usepackage[T1]{fontenc}
\usepackage{amsmath, amssymb, amsthm, mathtools}
\usepackage{geometry}
\usepackage{hyperref}
\usepackage{graphicx}
\usepackage{booktabs}
\usepackage{algorithm}
\usepackage{algpseudocode}
\usepackage{natbib}

\geometry{margin=1in}

\theoremstyle{plain}
\newtheorem{theorem}{Theorem}[section]
\newtheorem{lemma}[theorem]{Lemma}
\newtheorem{corollary}[theorem]{Corollary}

\theoremstyle{definition}
\newtheorem{definition}[theorem]{Definition}
\newtheorem{assumption}[theorem]{Assumption}

\hypersetup{colorlinks=true,linkcolor=cyan,citecolor=cyan,urlcolor=cyan}


\usepackage[top=1in, bottom=1in, left=1in, right=1in]{geometry}
\usepackage{draftwatermark}
\SetWatermarkText{SOVEREIGN PROTOCOL // ID: D1019B9}
\SetWatermarkScale{0.5}
\SetWatermarkLightness{0.94}
\begin{document}

\title{\textbf{
  Isomorphic Cryptographic Tethering: Using Physical Objects as Thermodynamic Memory Locks}}
\author{Rafael D. De Paz \\ \textit{Sovereign Architect} \\ \texttt{me@rdepaz.com}}
\date{2026-03-23}
\maketitle

\begin{abstract}
Advancing beyond traditional Hardware Security Modules (HSMs) and Physically Unclonable Functions (PUFs), this paper explores Isomorphic Cryptographic Tethering. We formalize how the literal thermodynamic and geometric topologies of specific physical objects (e.g., structured crystalline lattices or specialized jewelry) can be mapped identically to the hash functions of an access protocol. This creates a hard physical tether to reality, bounding digital information systems within strict thermodynamic constraints to prevent total systemic amnesia.
\end{abstract}

\vspace{1em}
\noindent\textbf{Keywords:} Vanguard Architecture, Systems Engineering, Canonical Optimization, Topological Security.

\vspace{0.5em}
\noindent\textbf{2026 MSC:} 68M14, 68Q85, 94A60, 81P68.

% =============================================================================
% 1. INTRODUCTION
% =============================================================================
\section{Introduction}
These foundational principles operate on established invariants \cite{suh2007physical,pappu2002physical}.\label{sec:intro}
This paper formulates the mathematical boundaries required to prove the stated thesis. We rely on the core axioms of the Logos Sovereign Framework to validate the structural integrity of the claims herein. Within modern macroscopic systems, unbounded entropy creates catastrophic localized failures; thus, rigorous constraints must be applied to network and execution topology.

% =============================================================================
% 2. PRELIMINARIES
% =============================================================================
\section{Preliminaries}\label{sec:prelim}
\subsection{Axiomatic Foundations}
Let $\mathcal{S}$ define the state boundaries of the substrate macro-node. The overarching theorem dictates that absolute minimization of operational latency maps isomorphically to maximum truth retention...

% =============================================================================
% 3. SYSTEMIC ARCHITECTURE
% =============================================================================
\section{Systemic Architecture}\label{sec:architecture}
We define the exact topological constraints...
\begin{theorem}[State Conservation]\label{thm:conservation}
The integration of a constraint matrix directly enforces state stability over $\Delta t$.
\end{theorem}
\begin{proof}
By formal application of Landauer's principle and structural thermodynamics, the bounds hold true...
\end{proof}

% =============================================================================
% 4. CONCLUSION
% =============================================================================
\section{Conclusion}\label{sec:conclusion}
The logical constraints of the system necessitate the conclusions drawn in the abstract, reinforcing the overarching theorem of thermodynamic alignment and computational sovereignty.

% =============================================================================
% REFERENCES
% =============================================================================

\vspace{2em}
\noindent\hrulefill
\vspace{1em}



% =============================================================================

% =============================================================================
% References & Cryptographic Lineage
% =============================================================================
\bibliographystyle{plainnat}
\bibliography{references}

\vspace{3em}
\noindent\hrulefill
\vspace{1em}

\end{document}

\end{document}